% КУРСОВАЯ РАБОТА
% Тема: "Социальные проекты для молодежи"
% Оформление по ГОСТ 7.32-2017 и требованиям нормоконтроля
% Компиляция: pdfLaTeX (Overleaf)

\documentclass[14pt,a4paper]{extarticle}

% ПОДКЛЮЧЕНИЕ ПАКЕТОВ ДЛЯ pdfLaTeX
\usepackage[T2A]{fontenc}                    % Кодировка шрифтов
\usepackage[utf8]{inputenc}                  % Кодировка исходного текста
\usepackage[english,russian]{babel}          % Локализация и переносы

% Шрифт Times New Roman (аналог для pdfLaTeX)
\usepackage{tempora}

% НАСТРОЙКА ПОЛЕЙ СТРАНИЦЫ (ГОСТ 7.32-2017)
% Левое - 30 мм, Правое - 15 мм, Верхнее - 20 мм, Нижнее - 20 мм
\usepackage[
    left=30mm,
    right=15mm,
    top=20mm,
    bottom=20mm
]{geometry}

% МЕЖСТРОЧНЫЙ ИНТЕРВАЛ 1.5 (ГОСТ)
\usepackage{setspace}
\onehalfspacing

% АБЗАЦНЫЙ ОТСТУП 1.25 см (ГОСТ)
\usepackage{indentfirst}                     % Отступ первого абзаца
\setlength{\parindent}{1.25cm}               % Размер отступа

% НАСТРОЙКА ЗАГОЛОВКОВ (ГОСТ)
\usepackage{titlesec}

% Формат заголовка раздела (section): прописные буквы, по центру, жирный
\titleformat{\section}
    {\normalfont\bfseries\centering}
    {\thesection}
    {1em}
    {\MakeUppercase}

% Формат заголовка подраздела (subsection)
\titleformat{\subsection}
    {\normalfont\bfseries}
    {\thesubsection}
    {1em}
    {}

% Интервалы до и после заголовков
\titlespacing*{\section}
    {0pt}{2ex plus 1ex minus .2ex}{1.5ex plus .2ex}
\titlespacing*{\subsection}
    {\parindent}{2ex plus 1ex minus .2ex}{1ex plus .2ex}

% НУМЕРАЦИЯ СТРАНИЦ (ГОСТ) - внизу по центру
\usepackage{fancyhdr}
\pagestyle{fancy}
\fancyhf{}
\fancyfoot[C]{\thepage}
\renewcommand{\headrulewidth}{0pt}
\renewcommand{\footrulewidth}{0pt}

% Для страниц с началом разделов
\fancypagestyle{plain}{
    \fancyhf{}
    \fancyfoot[C]{\thepage}
    \renewcommand{\headrulewidth}{0pt}
    \renewcommand{\footrulewidth}{0pt}
}

% НАСТРОЙКА ОГЛАВЛЕНИЯ (ГОСТ)
\usepackage{tocloft}

% Название "СОДЕРЖАНИЕ" по центру, жирным
\renewcommand{\cfttoctitlefont}{\hfill\bfseries}
\renewcommand{\cftaftertoctitle}{\hfill}

% Убираем жирность в самом оглавлении
\renewcommand{\cftsecfont}{\normalfont}
\renewcommand{\cftsecpagefont}{\normalfont}

% Точки-заполнители для всех уровней
\renewcommand{\cftsecleader}{\cftdotfill{\cftdotsep}}

% Отступы в оглавлении
\cftsetindents{section}{0em}{1.5em}
\cftsetindents{subsection}{1.5em}{2.3em}

% НАСТРОЙКА СПИСКОВ
\usepackage{enumitem}
\setlist{nosep, leftmargin=\parindent}

% ГИПЕРССЫЛКИ (опционально, для кликабельного оглавления)
\usepackage{hyperref}
\hypersetup{
    colorlinks=false,
    pdfborder={0 0 0},
}

% НАЧАЛО ДОКУМЕНТА
\begin{document}

% НУМЕРАЦИЯ НАЧИНАЕТСЯ СО 2-Й СТРАНИЦЫ (1-я - титульник)
\setcounter{page}{2}

% СОДЕРЖАНИЕ
\renewcommand{\contentsname}{СОДЕРЖАНИЕ}
\tableofcontents
\thispagestyle{fancy}

% ВВЕДЕНИЕ
\newpage
\section*{ВВЕДЕНИЕ}
\addcontentsline{toc}{section}{ВВЕДЕНИЕ}

\textbf{Актуальность исследования.} Молодежь является стратегическим ресурсом развития любого государства, от социальной позиции которой зависит решение вопросов народосбережения, достижения социально-экономического прогресса и повышения благосостояния общества. В современных социально-политических условиях центр тяжести социальной политики государства закономерно смещается в сторону молодежи как социально-демографической группы, численность которой в России составляет около 30 млн человек в возрасте от 14 до 35 лет.

Принятие Федерального закона <<О молодежной политике в Российской Федерации>> (2020) и утверждение <<Стратегии реализации молодежной политики в Российской Федерации на период до 2030 года>> (2024) свидетельствуют о системном подходе государства к работе с молодым поколением. Стратегическая цель молодежной политики определяется как <<воспитание гармонично развитого, высоконравственного и ответственного поколения российских граждан, способного обеспечить суверенитет, конкурентоспособность и дальнейшее развитие России>>.

В этом контексте социальное проектирование выступает одним из наиболее эффективных инструментов формирования социальной активности, гражданской идентичности и позитивной социализации молодежи. Социально-проектная деятельность позволяет молодым людям применять базовые знания и умения для решения реальных социальных проблем, развивать самостоятельность, критическое мышление, коммуникативные и организаторские способности.

Федеральные проекты <<Социальная активность>> и <<Молодежь России>> направлены на создание условий для развития добровольчества как ключевого элемента социальной ответственности гражданского общества. Программа <<Росмолодежь.Гранты>> предоставляет молодым гражданам возможность получить финансовую поддержку на реализацию социально значимых проектов по 18 номинациям, включая экологическое просвещение, развитие добровольчества и патриотическое воспитание.

Вместе с тем существуют противоречия между наличием общественного запроса на социально активную, адаптированную личность и недостаточной систематизацией работы по вовлечению молодежи в социальное проектирование; между востребованностью социально направленных практик и ограниченной информированностью молодежи о возможностях участия в них. Это обусловливает актуальность исследования социальных проектов как механизма развития молодежного потенциала.

\textbf{Степень разработанности проблемы.} Теоретические основы социального проектирования рассмотрены в трудах О.\,Ю.~Алюковой, Ж.\,Т.~Тощенко, А.\,А.~Сафиной, И.\,Т.~Тазетдиновой. Социально-проектная деятельность молодежи как инструмент социализации исследуется Л.\,С.~Пастуховой, С.\,В.~Ивановой, Л.\,В.~Байбородовой, Л.\,Н.~Серебренниковым. Проблемы волонтерства и добровольческой активности молодежи анализируются М.\,В.~Певной, А.\,Н.~Тарасовой, Д.\,Ф.~Телепаевой. Вопросы гражданско-патриотического воспитания посредством проектной деятельности изучают Т.\,Г.~Деревягина, Е.\,В.~Деева, Г.\,Я.~Гревцева, М.\,В.~Афанасьев. Социологические аспекты молодежной политики представлены в работах О.\,А.~Аникеевой, Е.\,А.~Воронцовой, В.\,Т.~Лисовского.

\textbf{Объект исследования} "--- социально-проектная деятельность в системе молодежной политики Российской Федерации.

\textbf{Предмет исследования} "--- социальные проекты как механизм развития социальной активности, гражданской идентичности и позитивной социализации молодежи.

\textbf{Цель исследования} "--- комплексный анализ теоретических основ, направлений и механизмов реализации социальных проектов для молодежи в современной России.

\textbf{Задачи исследования:}
\begin{enumerate}[leftmargin=\parindent, labelsep=0.5em, itemsep=0pt, parsep=0pt]
    \item Раскрыть понятие и сущность социального проектирования в контексте работы с молодежью.
    \item Систематизировать виды и направления социальных проектов для молодежи.
    \item Проанализировать роль социальных проектов в социализации молодежи и формировании гражданской позиции.
    \item Исследовать механизмы грантовой поддержки молодежных социальных инициатив.
    \item Определить критерии оценки эффективности социальных проектов и выявить проблемы вовлечения молодежи в проектную деятельность.
\end{enumerate}

\textbf{Методы исследования:} теоретический анализ научной литературы и нормативно-правовых документов, систематизация и обобщение, сравнительный анализ, метод описания и интерпретации.

\textbf{Теоретическая значимость} исследования состоит в обобщении научных подходов к пониманию социального проектирования как социально-педагогической технологии, направленной на позитивную социализацию молодежи, а также в уточнении содержания понятия <<социальная активность молодежи>> в контексте проектной деятельности.

\textbf{Практическая значимость} работы определяется возможностью использования её результатов при разработке программ воспитательной работы в образовательных организациях, проектировании молодежных добровольческих инициатив, а также в деятельности некоммерческих организаций, реализующих социальные проекты.

\textbf{Структура работы.} Курсовая работа состоит из введения, трёх глав, заключения и списка использованных источников. В первой главе рассматриваются теоретические основы социального проектирования в работе с молодежью. Вторая глава посвящена анализу роли социальных проектов в развитии молодежи. В третьей главе исследуются механизмы поддержки и оценка эффективности социальных проектов.

% ГЛАВА 1
\newpage
\section{ТЕОРЕТИЧЕСКИЕ ОСНОВЫ СОЦИАЛЬНОГО ПРОЕКТИРОВАНИЯ В РАБОТЕ С МОЛОДЕЖЬЮ}

\subsection{Понятие и сущность социального проектирования}

% Здесь будет текст параграфа 1.1

\subsection{Виды и направления социальных проектов для молодежи}

% Здесь будет текст параграфа 1.2

\subsection{Нормативно-правовая база молодежной политики в России}

% Здесь будет текст параграфа 1.3

% ГЛАВА 2
\newpage
\section{РОЛЬ СОЦИАЛЬНЫХ ПРОЕКТОВ В РАЗВИТИИ МОЛОДЕЖИ}

\subsection{Социализация молодежи через проектную деятельность}

% Здесь будет текст параграфа 2.1

\subsection{Волонтерство как форма социального проектирования}

% Здесь будет текст параграфа 2.2

\subsection{Гражданско-патриотическое воспитание молодежи посредством социальных проектов}

% Здесь будет текст параграфа 2.3

% ГЛАВА 3
\newpage
\section{МЕХАНИЗМЫ ПОДДЕРЖКИ И ОЦЕНКА ЭФФЕКТИВНОСТИ СОЦИАЛЬНЫХ ПРОЕКТОВ}

\subsection{Грантовая поддержка молодежных социальных инициатив}

% Здесь будет текст параграфа 3.1

\subsection{Критерии оценки эффективности социальных проектов}

% Здесь будет текст параграфа 3.2

\subsection{Проблемы и перспективы развития социального проектирования молодежи}

% Здесь будет текст параграфа 3.3

% ЗАКЛЮЧЕНИЕ
\newpage
\section*{ЗАКЛЮЧЕНИЕ}
\addcontentsline{toc}{section}{ЗАКЛЮЧЕНИЕ}

% Здесь будет текст заключения

% СПИСОК ИСПОЛЬЗОВАННЫХ ИСТОЧНИКОВ
\newpage
\section*{СПИСОК ИСПОЛЬЗОВАННЫХ ИСТОЧНИКОВ}
\addcontentsline{toc}{section}{СПИСОК ИСПОЛЬЗОВАННЫХ ИСТОЧНИКОВ}

% Здесь будет список литературы

\end{document}
